\documentclass[10pt,a4paper]{amsart}
\usepackage[margin=19mm]{geometry}
\usepackage{amsmath,amssymb,mathtools}
\usepackage[scr=rsfs,cal=euler]{mathalfa}
\usepackage{xfrac}
\usepackage[unicode,hidelinks,bookmarks=false]{hyperref}
\newcommand{\ie}{i.e.\ }
\newcommand{\cf}{cf.\ }
\newcommand{\resp}{respectively\ }
\newcommand{\Subsection}{\subsection}
\providecommand{\nobreakdash}{\nobreak\hbox{-}}
\setlength{\emergencystretch}{2em}
\multlinegap0pt
\usepackage{bm}
\usepackage{bbm}
\usepackage{dsfont}
\usepackage{xspace}
\usepackage{cancel}
\usepackage{mathtools}
\usepackage{amsthm}
\makeatletter
\@ifundefined{equation}{\newtheorem{equation}{Equation}}{}
\@ifundefined{lemma}{\newtheorem{lemma}[equation]{Lemma}}{}
\makeatother
\makeatletter
\newcommand{\C}{{\mathbb{C}}}
\renewcommand{\P}{{\mathbb{P}}}
\expandafter\ifx\csname pf\endcsname\relax
\newenvironment{pf}{\par{\it Proof.}}{{\it Q.E.D.}\par\vskip3pt}
\fi
\makeatother
\title[Meromorphic mappings into projective varieties intersecting ]{Meromorphic mappings into projective varieties intersecting arbitrary families of moving hypersurfaces}
\author{Si Duc Quang, Nguyen Linh Chi}
\date{Source: arXiv:2605.26034v1 (CC BY 4.0)}
\begin{document}
\maketitle

\noindent\emph{Source and licence.} Meromorphic mappings into projective varieties intersecting arbitrary families of moving hypersurfaces. Version arXiv:2605.26034v1 (CC BY 4.0), distributed under the Creative Commons Attribution 4.0 International licence, as recorded in the arXiv licence metadata for this exact version and shown on the abstract page https://arxiv.org/abs/2605.26034v1. Full rights notice: licenses/arxiv-2605.26034-CC-BY-4.0.txt.\par\smallskip

\noindent\textit{Editorial scope.} Counted unit: the Lemma whose source label is \texttt{4.1} in the article (source0.tex lines 595--597, source label \texttt{4.1}), delivered together with the article's own complete original proof of it (source0.tex lines 599--649, one \texttt{proof} environment of 51 lines). No mathematics is written, repaired or re-derived: statement and proof are the article's own text line for line. Required context retained uncounted: none. Every definition, display and label of those lines is retained unchanged. Omitted from this package and not redistributed: 1--594, 598--598, 650--788. Nothing inside a retained excerpt is omitted or altered: every retained body is delivered exactly as the article's lines are written, so no omission receipt of any kind accompanies this package. Citations: the retained excerpts carry no citation command at all, so no bibliography entry of the article is transcribed and no reference is redistributed. Reference adaptation: none was needed and none was made; every reference inside the delivered unit resolves inside it, so no unresolved reference or citation is produced. Printed slips kept literal: no printed word, symbol, hypothesis or number of the counted statement or of its proof was altered. Layout and class: the article's own class is \texttt{amsart} with packages including amsmath, amssymb; the wrapper is the lane's pinned \texttt{amsart} editorial wrapper at 10pt on A4, which restates the theorem-like environments the retained text needs and adds the article's own macro definitions that the retained text needs (\texttt{\textbackslash C}, \texttt{\textbackslash P}), with extra wrapper packages (bm, bbm, dsfont, xspace, cancel, mathtools). The delivered numbering is the unit's own. Assets: no retained span contains a figure environment, a graphics-inclusion command, a PDF-page inclusion, a table or a TikZ picture; no image file is copied into this package, the package declares no asset dependency and no image file is redistributed with it. Licence: version arXiv:2605.26034v1 of this article is distributed under the Creative Commons Attribution 4.0 International licence, as recorded in the arXiv licence metadata for this exact version and shown on the abstract page https://arxiv.org/abs/2605.26034v1. A full rights notice is delivered at licenses/arxiv-2605.26034-CC-BY-4.0.txt. Characters: every retained excerpt is pure ASCII, so the delivered engine, XeLaTeX with its default Latin Modern text font, reports no missing character.
\begin{lemma}\label{4.1}
Let $f$ be a nonconstant meromorphic mappings from $\C^m$ into a projective subvariety $V\subset\P^n(\C)$ of dimension $k$. Let $\mathcal{Q} = \{Q_1, \ldots, Q_q\}$ be a family of moving hypersurfaces in $\P^n(\C)$ which has distributive constant $\Delta_{\mathcal Q,V}$ with respect to $V\ (N> k)$ and satisfying the B\'ezout property with respect to $f$. Let $n_f$ be the algebraic dimension of $f$ over the field $\mathcal K_{\mathcal Q}$. Then there exists a subset $\Gamma \subset \{1,\ldots,q\}$ such that $\sharp\Gamma\ge q -\Delta_{\mathcal Q,V}\left(k-\left\lceil\frac{n_f}{2}\right\rceil\right)$ and the family $\mathcal P= \{Q_i : i \in \Gamma\}$ satisfies $\Delta_{\mathcal P,f}\le \frac{2k-n_f+1}{n_f+1}\Delta_{\mathcal Q,V}$.
\end{lemma}

\begin{proof} 
Let $\{R_1,\ldots,R_\lambda\}$ a minimal subset of $\mathcal K_{\mathcal Q}[x_0,\ldots,x_n]$, which generates $I(f(\C^m))$. 
For $z$ is not a pole of any $R_i\ (1\le i\le \lambda)$, set $V_z=\bigcap_{i=1}^\lambda R_i(z)^*$. Then $\dim V_z=n_f$ for generic points $z$.

For any subset $\Gamma \subset \{1,\ldots,q\}$, denote
$$Q_\Gamma(z) = \bigcap_{i \in \Gamma}Q_i(z)^* \cap V_z, \quad c(\Gamma) = \operatorname{codim}_{V_z} Q_\Gamma(z)$$
for generic points $z$. Note that
$$\Delta_{\mathcal Q,f}= \max_{\Gamma \ne \emptyset} \frac{\sharp\Gamma}{c(\Gamma)}.$$
If $\Delta_{\mathcal Q,f}\le \frac{2k-n_f+1}{n_f+1}\Delta_{\mathcal Q,V}$, then the conclusion holds with $\Gamma = \{1,\ldots,q\}$.

Assume $\Delta_{\mathcal Q,f}>\frac{2k-n_f+1}{n_f+1}\Delta_{\mathcal Q,V}$. Let
$$\mathcal{S} = \left\{ \Gamma \subset \{1,\ldots,q\} : \frac{\sharp\Gamma}{c(\Gamma)} >\frac{2k-n_f+1}{n_f+1}\right\}.$$
Then $\mathcal{S}$ is nonempty. Choose $\Gamma_1 \in \mathcal{S}$ such that $\sharp\Gamma_1$ is maximal. Since 
\begin{align*}
\sharp\Gamma_1&\le \Delta_{\mathcal Q,V}\left(\dim V-\dim V\cap\bigl(\bigcap_{j\in\Gamma_1}Q_j(z)^*\bigl)\right)\\
&\le \Delta_{\mathcal Q,V}\left(\dim V-\dim V_z\cap\bigl(\bigcap_{j\in\Gamma_1}Q_j(z)^*\bigl)\right)\\
&=\Delta_{\mathcal Q,V}\left(k-n_f+c(\Gamma_1)\right),
\end{align*} 
we have
$$\frac{2k-n_f+1}{n_f+1}\Delta_{\mathcal Q,V}<\frac{\sharp\Gamma_1}{c(\Gamma_1)}\le \frac{\left(k-n_f+c(\Gamma_1)\right)}{c(\Gamma_1)}\Delta_{\mathcal Q,V}.$$
This implies that $c(\Gamma_1)<\frac{n_f+1}{2}$ and then $c(\Gamma_1)\le \left\lfloor\frac{n_f}{2}\right\rfloor$.

\medskip
\noindent \textbf{Claim.} For any $\Gamma_2 \in \mathcal{S}\setminus\{\Gamma_1\}$, we have $\Gamma_1 \cap \Gamma_2 \ne \emptyset$.
\medskip

\noindent \textit{Proof of Claim.}
Assume $\Gamma_1 \cap \Gamma_2 = \emptyset$. Similar as above, one has $c(\Gamma_2)<\frac{n_f+1}{2}$.  Using the B\'ezout property, we have
$$c(\Gamma_1 \cup \Gamma_2) \le c(\Gamma_1) + c(\Gamma_2)<n_f+1.$$
This yields that $\bigcap_{j\in \Gamma_1 \cup \Gamma_2}(Q_j(z)^*\cap V_z)\ne\varnothing$ for generic points $z$.

Also, from
$$\frac{\sharp\Gamma_1}{c(\Gamma_1)} > \frac{2k-n_f+1}{n_f+1}\Delta_{\mathcal Q,V} \quad\text{ and }\quad \frac{\sharp\Gamma_2}{c(\Gamma_2)} > \frac{2k-n_f+1}{n_f+1}\Delta_{\mathcal Q,V},$$
by using the B\'ezout property we obtain
$$\frac{\sharp\Gamma_1 + \sharp\Gamma_2}{c(\Gamma_1 \cup \Gamma_2)}\ge \frac{\sharp\Gamma_1+\sharp\Gamma_2}{c(\Gamma_1) + c(\Gamma_2)} > \frac{2k-n_f+1}{n_f+1}\Delta_{\mathcal Q,V}.$$
Thus $\Gamma_1 \cup \Gamma_2 \in \mathcal{S}$, contradicting the maximality of $\sharp\Gamma_1$. 
The claim is proved.

Now set
$$\Gamma = \{1,\ldots,q\} \setminus \Gamma_1, \quad \mathcal{P} = \{Q_i : i \in \Gamma\}.$$
Then
$$\sharp\Gamma\ge q -\Delta_{\mathcal Q,V}\left(k-n_f+c(\Gamma_1)\right)\ge q -\Delta_{\mathcal Q,V}\left(k-\left\lceil\frac{n_f}{2}\right\rceil\right).$$
Finally, suppose there exists $\Gamma' \subset \Gamma$ such that
$$
\frac{|\Gamma'|}{c(\Gamma')} > \frac{2k-n_f+1}{n_f+1}\Delta_{\mathcal Q,V}.
$$
Then $\Gamma' \in \mathcal{S}$ but $\Gamma' \cap \Gamma_1 = \emptyset$, contradicting the claim. 
Hence
$$\Delta_{\mathcal P,V}\le \frac{2k-n_f+1}{n_f+1}\Delta_{\mathcal Q,V}.$$
This completes the proof.
\end{proof}

\end{document}
